The rise of generative artificial intelligence in education presents a paradox, these tools can solve in seconds the very programming exercises students need to learn to solve on their own. Banning them is counterproductive in a context where their professional use is already inevitable, and allowing them without control leads to passive learning and cognitive dependency.

This Bachelor's Thesis presents \textbf{SocratiCode}, a web-based programming tutoring platform that proposes a third way. Rather than generating solutions, the system uses a locally-hosted language model instructed through prompt engineering to act as a Socratic tutor. It guides students with targeted questions that lead them to discover errors and solutions on their own, and explicitly refuses to provide corrected code under any circumstances.

The platform integrates a code editor with sandboxed execution, a real-time communication channel based on \textit{Server-Sent Events}, and a dual moderation system that safeguards the educational environment. By running entirely on local infrastructure, it requires no commercial API dependencies and exposes no academic data to third parties.

The interface, designed as a single-page application, brings together the code editor, the output terminal, and the chat with the tutor on a single screen. This layout allows students to see at a glance the error they made, the output it produced, and the Socratic hint the tutor is offering, eliminating the need to switch between applications during the learning process.

\keywords{Large language models (LLM), prompt engineering, generative artificial intelligence, socratic tutor, programming education, secure code execution}
